%%%%%%%%%%%%%%%%%%%%%%%%%%%%%%%%%%%%%%%%%%%%%%%%%%%%%%%%%%%%%%%%%%%%%%
% Lớp 9 | Chương VI. Hàm số y = ax² (a ≠ 0). Phương trình bậc hai một ẩn | Định lí Viète và ứng dụng
% Dạng 9D6?20-2: Tính giá trị biểu thức đối xứng của hai nghiệm
%%%%%%%%%%%%%%%%%%%%%%%%%%%%%%%%%%%%%%%%%%%%%%%%%%%%%%%%%%%%%%%%%%%%%%


% nguồn: 2026-11-14_09_de_tuyen_sinh_10_tphcm_2025_2026 (nhập 14/11/2026)
\begin{bt}%[9D6H20-2]%[ma:df44b7b6]
% nguon: Đề thi chính thức Tuyển sinh 10 TP.HCM năm học 2025-2026 – Bài 2
Cho phương trình $2x^2 - 7x + 4 = 0$.
\begin{enumerate}[a)]
	\item Chứng minh phương trình trên có hai nghiệm phân biệt $x_1, x_2$.
	\item Không giải phương trình, hãy tính giá trị của biểu thức $A = x_1(3x_2 + x_1) + x_2^2$.
\end{enumerate}
\loigiai{
\begin{enumerate}[a)]
	\item Xét phương trình $2x^2 - 7x + 4 = 0$, ta có:
	$\Delta = b^2 - 4ac = (-7)^2 - 4\cdot 2\cdot 4 = 49 - 32 = 17 > 0$.\\
	Vì $\Delta > 0$ nên phương trình luôn có hai nghiệm phân biệt $x_1, x_2$.
	\item Theo hệ thức Viète, ta có:
	$\begin{cases} x_1 + x_2 = -\dfrac{b}{a} = \dfrac{7}{2} \\ x_1 x_2 = \dfrac{c}{a} = \dfrac{4}{2} = 2. \end{cases}$\\
	Biến đổi biểu thức $A$:
	$A = x_1(3x_2 + x_1) + x_2^2 = 3x_1 x_2 + x_1^2 + x_2^2 = 3x_1 x_2 + (x_1 + x_2)^2 - 2x_1 x_2 = (x_1 + x_2)^2 + x_1 x_2$.\\
	Thay $x_1 + x_2 = \dfrac{7}{2}$ và $x_1 x_2 = 2$ vào biểu thức $A$, ta được:
	$A = \left(\dfrac{7}{2}\right)^2 + 2 = \dfrac{49}{4} + 2 = \dfrac{57}{4}$.\\
	Vậy $A = \dfrac{57}{4}$.
\end{enumerate}
}
\end{bt}

% nguồn: 2026-11-14_09_de_minh_hoa_tuyen_sinh_10_tphcm (nhập 14/11/2026)
\begin{bt}%[9D6H20-2]%[ma:8a99860b]
% nguon: Đề thi tham khảo Tuyển sinh 10 TP.HCM – Phụ lục Sở GD&ĐT TP.HCM – Bài 2
Cho phương trình $2x^2 - 5x + 1 = 0$.
\begin{enumerate}[a)]
	\item Chứng minh phương trình trên có hai nghiệm phân biệt.
	\item Không giải phương trình, hãy tính giá trị của biểu thức $A = x_1(x_1 + 2024) + x_2(x_2 + 2025) - x_2$.
\end{enumerate}
\loigiai{
\begin{enumerate}[a)]
	\item Xét phương trình $2x^2 - 5x + 1 = 0$, ta có:
	$\Delta = b^2 - 4ac = (-5)^2 - 4\cdot 2\cdot 1 = 25 - 8 = 17 > 0$.\\
	Vì $\Delta > 0$ nên phương trình luôn có hai nghiệm phân biệt $x_1, x_2$.
	\item Theo định lí Viète, ta có:
	$\begin{cases} x_1 + x_2 = -\dfrac{b}{a} = \dfrac{5}{2} \\ x_1 x_2 = \dfrac{c}{a} = \dfrac{1}{2}. \end{cases}$\\
	Biến đổi biểu thức $A$:
	$A = x_1^2 + 2024x_1 + x_2^2 + 2025x_2 - x_2 = (x_1^2 + x_2^2) + 2024(x_1 + x_2) = (x_1 + x_2)^2 - 2x_1 x_2 + 2024(x_1 + x_2)$.\\
	Thay $x_1 + x_2 = \dfrac{5}{2}$ và $x_1 x_2 = \dfrac{1}{2}$ vào biểu thức $A$, ta được:
	$A = \left(\dfrac{5}{2}\right)^2 - 2\cdot\dfrac{1}{2} + 2024\cdot\dfrac{5}{2} = \dfrac{25}{4} - 1 + 5060 = \dfrac{21}{4} + 5060 = \dfrac{20261}{4} = 5065{,}25$.\\
	Vậy $A = \dfrac{20261}{4} = 5065{,}25$.
\end{enumerate}
}
\end{bt}

% nguồn: 2026-11-15_09_de_tham_khao_quan_1_de_1_va_de_2 (nhập 15/11/2026)
\begin{bt}%[9D6H20-2]%[ma:ae356dc7]
% nguon: Đề tham khảo Tuyển sinh 10 TP.HCM 2025-2026 – Quận 1 – Đề 1 – Bài 2
Cho phương trình $2x^2 - 9x + 3 = 0$.
\begin{enumerate}[a)]
	\item Chứng tỏ phương trình trên có hai nghiệm phân biệt.
	\item Gọi $x_1, x_2$ là hai nghiệm của phương trình. Tính giá trị của biểu thức:
	$M = x_1(2025 - x_1) + x_2(2026 - x_2) - x_2$.
\end{enumerate}
\loigiai{
\begin{enumerate}[a)]
	\item Xét phương trình $2x^2 - 9x + 3 = 0$, ta có:
	$\Delta = b^2 - 4ac = (-9)^2 - 4\cdot 2\cdot 3 = 81 - 24 = 57 > 0$.\\
	Vì $\Delta > 0$ nên phương trình luôn có hai nghiệm phân biệt $x_1, x_2$.
	\item Theo định lí Viète, ta có:
	$\begin{cases} x_1 + x_2 = -\dfrac{b}{a} = \dfrac{9}{2} \\ x_1 x_2 = \dfrac{c}{a} = \dfrac{3}{2}. \end{cases}$\\
	Biến đổi biểu thức $M$:
	$M = 2025x_1 - x_1^2 + 2026x_2 - x_2^2 - x_2 = 2025(x_1 + x_2) - (x_1^2 + x_2^2) = 2025(x_1 + x_2) - \left[(x_1 + x_2)^2 - 2x_1 x_2\right]$.\\
	Thay $x_1 + x_2 = \dfrac{9}{2}$ và $x_1 x_2 = \dfrac{3}{2}$ vào biểu thức $M$, ta được:
	$M = 2025\cdot\dfrac{9}{2} - \left[\left(\dfrac{9}{2}\right)^2 - 2\cdot\dfrac{3}{2}\right] = \dfrac{18225}{2} - \left(\dfrac{81}{4} - 3\right) = \dfrac{18225}{2} - \dfrac{69}{4} = \dfrac{36381}{4} = 9095{,}25$.\\
	Vậy $M = \dfrac{36381}{4} = 9095{,}25$.
\end{enumerate}
}
\end{bt}

% nguồn: 2026-11-15_09_de_tham_khao_quan_1_de_1_va_de_2 (nhập 15/11/2026)
\begin{bt}%[9D6H20-2]%[ma:9bd51aa8]
% nguon: Đề tham khảo Tuyển sinh 10 TP.HCM 2025-2026 – Quận 1 – Đề 2 – Bài 2
Cho phương trình $3x^2 - 2x - 2 = 0$.
\begin{enumerate}[a)]
	\item Chứng tỏ phương trình trên có hai nghiệm phân biệt.
	\item Không giải phương trình, hãy tính giá trị của biểu thức $T = \dfrac{x_1}{x_2 - 1} + \dfrac{x_2}{x_1 - 1}$.
\end{enumerate}
\loigiai{
\begin{enumerate}[a)]
	\item Xét phương trình $3x^2 - 2x - 2 = 0$, ta có:
	$a\cdot c = 3\cdot(-2) = -6 < 0$.\\
	Vì tích $ac < 0$ nên phương trình luôn có hai nghiệm phân biệt trái dấu $x_1, x_2$.
	\item Theo định lí Viète, ta có:
	$\begin{cases} x_1 + x_2 = -\dfrac{b}{a} = \dfrac{2}{3} \\ x_1 x_2 = \dfrac{c}{a} = -\dfrac{2}{3}. \end{cases}$\\
	Quy đồng mẫu số biểu thức $T$:
	$T = \dfrac{x_1(x_1 - 1) + x_2(x_2 - 1)}{(x_2 - 1)(x_1 - 1)} = \dfrac{x_1^2 + x_2^2 - (x_1 + x_2)}{x_1 x_2 - (x_1 + x_2) + 1} = \dfrac{(x_1 + x_2)^2 - 2x_1 x_2 - (x_1 + x_2)}{x_1 x_2 - (x_1 + x_2) + 1}$.\\
	Thay $x_1 + x_2 = \dfrac{2}{3}$ và $x_1 x_2 = -\dfrac{2}{3}$ vào biểu thức $T$, ta được:
	$T = \dfrac{\left(\dfrac{2}{3}\right)^2 - 2\cdot\left(-\dfrac{2}{3}\right) - \dfrac{2}{3}}{-\dfrac{2}{3} - \dfrac{2}{3} + 1} = \dfrac{\dfrac{4}{9} + \dfrac{4}{3} - \dfrac{2}{3}}{-\dfrac{4}{3} + 1} = \dfrac{\dfrac{4}{9} + \dfrac{6}{9}}{-\dfrac{1}{3}} = \dfrac{\dfrac{10}{9}}{-\dfrac{1}{3}} = -\dfrac{10}{3}$.\\
	Vậy $T = -\dfrac{10}{3}$.
\end{enumerate}
}
\end{bt}

% nguồn: 2026-11-15_09_de_tham_khao_quan_3_de_1_va_de_2 (nhập 15/11/2026)
\begin{bt}%[9D6H20-2]%[ma:f5eddbd5]
% nguon: Đề tham khảo Tuyển sinh 10 TP.HCM 2025-2026 – Quận 3 – Đề 1 – Bài 2
Cho phương trình $2x^2 + 4x - 1 = 0$. Không giải phương trình:
\begin{enumerate}[a)]
	\item Chứng tỏ phương trình luôn có hai nghiệm phân biệt.
	\item Tính giá trị của biểu thức $A = \dfrac{1}{2}x_1(2x_1 + 3x_2) + \dfrac{1}{4}x_2(4x_2 + 6x_1)$.
\end{enumerate}
\loigiai{
\begin{enumerate}[a)]
	\item Xét phương trình $2x^2 + 4x - 1 = 0$, ta có hệ số $a = 2$ và $c = -1$.\\
	Vì tích $a\cdot c = 2\cdot(-1) = -2 < 0$ nên phương trình luôn có hai nghiệm phân biệt trái dấu $x_1, x_2$.
	\item Theo định lí Viète, ta có:
	$\begin{cases} x_1 + x_2 = -\dfrac{b}{a} = -\dfrac{4}{2} = -2 \\ x_1 x_2 = \dfrac{c}{a} = -\dfrac{1}{2}. \end{cases}$\\
	Biến đổi biểu thức $A$:
	$A = x_1^2 + \dfrac{3}{2}x_1 x_2 + x_2^2 + \dfrac{3}{2}x_1 x_2 = (x_1^2 + x_2^2) + 3x_1 x_2 = (x_1 + x_2)^2 + x_1 x_2$.\\
	Thay $x_1 + x_2 = -2$ và $x_1 x_2 = -\dfrac{1}{2}$ vào biểu thức $A$, ta được:
	$A = (-2)^2 + \left(-\dfrac{1}{2}\right) = 4 - \dfrac{1}{2} = \dfrac{7}{2} = 3{,}5$.\\
	Vậy $A = \dfrac{7}{2} = 3{,}5$.
\end{enumerate}
}
\end{bt}

% nguồn: 2026-11-15_09_de_tham_khao_quan_3_de_1_va_de_2 (nhập 15/11/2026)
\begin{bt}%[9D6H20-2]%[ma:30d20385]
% nguon: Đề tham khảo Tuyển sinh 10 TP.HCM 2025-2026 – Quận 3 – Đề 2 – Bài 2
Cho phương trình $3x^2 - 2x - 2 = 0$.
\begin{enumerate}[a)]
	\item Chứng tỏ phương trình luôn có hai nghiệm phân biệt $x_1, x_2$.
	\item Không giải phương trình, hãy tính giá trị của biểu thức $A = x_1^3 + x_2^3 - 3x_1^2 x_2 - 3x_1 x_2^2$.
\end{enumerate}
\loigiai{
\begin{enumerate}[a)]
	\item Xét phương trình $3x^2 - 2x - 2 = 0$, ta có:
	$a\cdot c = 3\cdot(-2) = -6 < 0$.\\
	Vì tích $ac < 0$ nên phương trình luôn có hai nghiệm phân biệt trái dấu $x_1, x_2$.
	\item Theo định lí Viète, ta có:
	$\begin{cases} x_1 + x_2 = -\dfrac{b}{a} = \dfrac{2}{3} \\ x_1 x_2 = \dfrac{c}{a} = -\dfrac{2}{3}. \end{cases}$\\
	Biến đổi biểu thức $A$:
	$A = (x_1^3 + x_2^3) - 3x_1 x_2 (x_1 + x_2) = (x_1 + x_2)(x_1^2 - x_1 x_2 + x_2^2) - 3x_1 x_2(x_1 + x_2) = (x_1 + x_2)\left[(x_1 + x_2)^2 - 3x_1 x_2\right] - 3x_1 x_2(x_1 + x_2) = (x_1 + x_2)\left[(x_1 + x_2)^2 - 6x_1 x_2\right]$.\\
	Thay $x_1 + x_2 = \dfrac{2}{3}$ và $x_1 x_2 = -\dfrac{2}{3}$ vào biểu thức $A$, ta được:
	$A = \dfrac{2}{3}\left[\left(\dfrac{2}{3}\right)^2 - 6\cdot\left(-\dfrac{2}{3}\right)\right] = \dfrac{2}{3}\left(\dfrac{4}{9} + 4\right) = \dfrac{2}{3} \cdot \dfrac{40}{9} = \dfrac{80}{27}$.\\
	Vậy $A = \dfrac{80}{27}$.
\end{enumerate}
}
\end{bt}

% nguồn: 2026-11-15_09_de_tham_khao_quan_4_de_1_va_de_2 (nhập 15/11/2026)
\begin{bt}%[9D6H20-2]%[ma:4215423d]
% nguon: Đề tham khảo Tuyển sinh 10 TP.HCM 2025-2026 – Quận 4 – Đề 1 – Bài 2
Cho phương trình $3x^2 - 7x + 3 = 0$.
\begin{enumerate}[a)]
	\item Chứng minh phương trình trên có hai nghiệm phân biệt.
	\item Không giải phương trình, gọi $x_1, x_2$ là hai nghiệm của phương trình đã cho, hãy tính giá trị của biểu thức:
	$A = \dfrac{18x_2}{5x_1 - x_2 + \dfrac{7}{3}} + \dfrac{18x_1}{5x_2 - x_1 + \dfrac{7}{3}}$.
\end{enumerate}
\loigiai{
\begin{enumerate}[a)]
	\item Xét phương trình $3x^2 - 7x + 3 = 0$, ta có:
	$\Delta = b^2 - 4ac = (-7)^2 - 4\cdot 3\cdot 3 = 49 - 36 = 13 > 0$.\\
	Vì $\Delta > 0$ nên phương trình luôn có hai nghiệm phân biệt $x_1, x_2$.
	\item Theo định lí Viète, ta có:
	$\begin{cases} x_1 + x_2 = -\dfrac{b}{a} = \dfrac{7}{3} \\ x_1 x_2 = \dfrac{c}{a} = \dfrac{3}{3} = 1. \end{cases}$\\
	Thay $\dfrac{7}{3} = x_1 + x_2$ vào mẫu thức của biểu thức $A$:
	\begin{itemize}
		\item Mẫu thức thứ nhất: $5x_1 - x_2 + \dfrac{7}{3} = 5x_1 - x_2 + (x_1 + x_2) = 6x_1$.
		\item Mẫu thức thứ hai: $5x_2 - x_1 + \dfrac{7}{3} = 5x_2 - x_1 + (x_1 + x_2) = 6x_2$.
	\end{itemize}
	Khi đó biểu thức $A$ trở thành:
	$A = \dfrac{18x_2}{6x_1} + \dfrac{18x_1}{6x_2} = 3\left(\dfrac{x_2}{x_1} + \dfrac{x_1}{x_2}\right) = 3\cdot\dfrac{x_1^2 + x_2^2}{x_1 x_2} = 3\cdot\dfrac{(x_1 + x_2)^2 - 2x_1 x_2}{x_1 x_2}$.\\
	Thay $x_1 + x_2 = \dfrac{7}{3}$ và $x_1 x_2 = 1$ vào, ta được:
	$A = 3\cdot\dfrac{\left(\dfrac{7}{3}\right)^2 - 2\cdot 1}{1} = 3\left(\dfrac{49}{9} - 2\right) = 3\cdot\dfrac{31}{9} = \dfrac{31}{3}$.\\
	Vậy $A = \dfrac{31}{3}$.
\end{enumerate}
}
\end{bt}

% nguồn: 2026-11-15_09_de_tham_khao_quan_4_de_1_va_de_2 (nhập 15/11/2026)
\begin{bt}%[9D6H20-2]%[ma:ff79bc1e]
% nguon: Đề tham khảo Tuyển sinh 10 TP.HCM 2025-2026 – Quận 4 – Đề 2 – Bài 2
Cho phương trình $3x^2 - 4x - 1 = 0$.
\begin{enumerate}[a)]
	\item Chứng tỏ phương trình trên có hai nghiệm phân biệt.
	\item Gọi $x_1, x_2$ là hai nghiệm của phương trình. Không giải phương trình, hãy tính giá trị của biểu thức $A = \dfrac{x_1}{x_2 - 1} + \dfrac{x_2}{x_1 - 1}$.
\end{enumerate}
\loigiai{
\begin{enumerate}[a)]
	\item Xét phương trình $3x^2 - 4x - 1 = 0$, ta có hệ số $a = 3$ và $c = -1$.\\
	Vì tích $a\cdot c = 3\cdot(-1) = -3 < 0$ nên phương trình luôn có hai nghiệm phân biệt trái dấu $x_1, x_2$.
	\item Theo định lí Viète, ta có:
	$\begin{cases} x_1 + x_2 = -\dfrac{b}{a} = \dfrac{4}{3} \\ x_1 x_2 = \dfrac{c}{a} = -\dfrac{1}{3}. \end{cases}$\\
	Quy đồng mẫu thức biểu thức $A$:
	$A = \dfrac{x_1(x_1 - 1) + x_2(x_2 - 1)}{(x_2 - 1)(x_1 - 1)} = \dfrac{x_1^2 + x_2^2 - (x_1 + x_2)}{x_1 x_2 - (x_1 + x_2) + 1} = \dfrac{(x_1 + x_2)^2 - 2x_1 x_2 - (x_1 + x_2)}{x_1 x_2 - (x_1 + x_2) + 1}$.\\
	Thay $x_1 + x_2 = \dfrac{4}{3}$ và $x_1 x_2 = -\dfrac{1}{3}$ vào biểu thức $A$, ta được:
	$A = \dfrac{\left(\dfrac{4}{3}\right)^2 - 2\cdot\left(-\dfrac{1}{3}\right) - \dfrac{4}{3}}{-\dfrac{1}{3} - \dfrac{4}{3} + 1} = \dfrac{\dfrac{16}{9} + \dfrac{2}{3} - \dfrac{4}{3}}{-\dfrac{5}{3} + 1} = \dfrac{\dfrac{16}{9} - \dfrac{6}{9}}{-\dfrac{2}{3}} = \dfrac{\dfrac{10}{9}}{-\dfrac{2}{3}} = -\dfrac{5}{3}$.\\
	Vậy $A = -\dfrac{5}{3}$.
\end{enumerate}
}
\end{bt}

% nguồn: 2026-11-15_09_de_tham_khao_quan_5_de_1_va_de_2 (nhập 15/11/2026)
\begin{bt}%[9D6H20-2]%[ma:52b068dc]
% nguon: Đề tham khảo Tuyển sinh 10 TP.HCM 2025-2026 – Quận 5 – Đề 1 – Bài 2
Cho phương trình $x^2 + 2x - 2025 = 0$.
\begin{enumerate}[a)]
    \item Chứng minh phương trình có hai nghiệm phân biệt trái dấu.
    \item Gọi $x_1$, $x_2$ là hai nghiệm của phương trình. Không giải phương trình, tính giá trị của biểu thức $A = \dfrac{(1 - x_1)x_2}{x_1} + \dfrac{(1 - x_2)x_1}{x_2}$.
\end{enumerate}
\loigiai{
\begin{enumerate}[a)]
    \item Phương trình $x^2 + 2x - 2025 = 0$ có hệ số $a = 1$ và $c = -2025$.\\
    Ta có tích $a \cdot c = 1 \cdot (-2025) = -2025 < 0$.\\
    Do $a$ và $c$ trái dấu nên phương trình luôn có hai nghiệm phân biệt trái dấu $x_1, x_2$.
    \item Theo định lí Viète, ta có:
    $\left\{\begin{aligned} &S = x_1 + x_2 = -2 \\ &P = x_1 x_2 = -2025. \end{aligned}\right.$\\
    Biến đổi biểu thức $A$:
    $\begin{aligned}
    A &= \dfrac{(1 - x_1)x_2}{x_1} + \dfrac{(1 - x_2)x_1}{x_2} = \dfrac{x_2^2(1 - x_1) + x_1^2(1 - x_2)}{x_1 x_2} \\
    &= \dfrac{(x_1^2 + x_2^2) - x_1 x_2(x_1 + x_2)}{x_1 x_2} = \dfrac{(x_1 + x_2)^2 - 2x_1 x_2 - x_1 x_2(x_1 + x_2)}{x_1 x_2} \\
    &= \dfrac{S^2 - 2P - P \cdot S}{P} = \dfrac{(-2)^2 - 2(-2025) - (-2025)(-2)}{-2025} \\
    &= \dfrac{4 + 4050 - 4050}{-2025} = -\dfrac{4}{2025}.
    \end{aligned}$\\
    Vậy $A = -\dfrac{4}{2025}$.
\end{enumerate}
}
\end{bt}

% nguồn: 2026-11-15_09_de_tham_khao_quan_5_de_1_va_de_2 (nhập 15/11/2026)
\begin{bt}%[9D6H20-2]%[ma:f94457b5]
% nguon: Đề tham khảo Tuyển sinh 10 TP.HCM 2025-2026 – Quận 5 – Đề 2 – Bài 2
Cho phương trình $x^2 - 4x - 6 = 0$.
\begin{enumerate}[a)]
    \item Chứng minh phương trình trên có hai nghiệm phân biệt $x_1$, $x_2$.
    \item Không giải phương trình, hãy tính giá trị của biểu thức: $A = \dfrac{x_1 x_2}{4 - x_1} + \dfrac{x_1 x_2}{4 - x_2}$.
\end{enumerate}
\loigiai{
\begin{enumerate}[a)]
    \item Phương trình $x^2 - 4x - 6 = 0$ có các hệ số $a = 1$ và $c = -6$.\\
    Vì $a \cdot c = 1 \cdot (-6) = -6 < 0$ nên phương trình luôn có hai nghiệm phân biệt $x_1, x_2$ trái dấu nhau.
    \item Theo định lí Viète, ta có:
    $\left\{\begin{aligned} &x_1 + x_2 = 4 \\ &x_1 x_2 = -6. \end{aligned}\right.$\\
    Từ $x_1 + x_2 = 4$, suy ra $4 - x_1 = x_2$ và $4 - x_2 = x_1$.\\
    Do đó biểu thức $A$ có thể viết lại thành:
    $A = \dfrac{x_1 x_2}{x_2} + \dfrac{x_1 x_2}{x_1} = x_1 + x_2 = 4$.\\
    Vậy $A = 4$.
\end{enumerate}
}
\end{bt}

% nguồn: 2026-11-15_09_de_tham_khao_quan_5_de_3_va_quan_6_de_1 (nhập 15/11/2026)
\begin{bt}%[9D6H20-2]%[ma:00644835]
% nguon: Đề tham khảo Tuyển sinh 10 TP.HCM 2025-2026 – Quận 5 – Đề 3 – Bài 2
Cho phương trình $4x^2 + 12x + 3 = 0$ \quad (1).
\begin{enumerate}[a)]
    \item Chứng minh phương trình (1) có hai nghiệm phân biệt.
    \item Gọi $x_1$, $x_2$ là hai nghiệm của phương trình (1). Không giải phương trình, tính giá trị của biểu thức $A = x_1^3 + x_2^3 + 2025$.
\end{enumerate}
\loigiai{
\begin{enumerate}[a)]
    \item Phương trình (1) có $a = 4$, $b' = 6$, $c = 3$.\\
    Biệt thức $\Delta' = (b')^2 - ac = 6^2 - 4 \cdot 3 = 36 - 12 = 24 > 0$.\\
    Vì $\Delta' > 0$ nên phương trình (1) luôn có hai nghiệm phân biệt $x_1, x_2$.
    \item Theo định lí Viète, ta có:
    $\left\{\begin{aligned} &S = x_1 + x_2 = -\dfrac{12}{4} = -3 \\ &P = x_1 x_2 = \dfrac{3}{4}. \end{aligned}\right.$\\
    Ta có hằng đẳng thức:
    $x_1^3 + x_2^3 = (x_1 + x_2)(x_1^2 - x_1 x_2 + x_2^2) = (x_1 + x_2)\left[(x_1 + x_2)^2 - 3x_1 x_2\right] = S(S^2 - 3P)$.\\
    Thay giá trị của $S$ và $P$ vào biểu thức $A$:
    $\begin{aligned}
    A &= S(S^2 - 3P) + 2025 = (-3)\left[(-3)^2 - 3 \cdot \dfrac{3}{4}\right] + 2025 \\
    &= (-3)\left(9 - \dfrac{9}{4}\right] + 2025 = (-3) \cdot \dfrac{27}{4} + 2025 = -\dfrac{81}{4} + 2025 = \dfrac{8019}{4} = 2004{,}75.
    \end{aligned}$\\
    Vậy $A = \dfrac{8019}{4} = 2004{,}75$.
\end{enumerate}
}
\end{bt}

% nguồn: 2026-11-15_09_de_tham_khao_quan_5_de_3_va_quan_6_de_1 (nhập 15/11/2026)
\begin{bt}%[9D6H20-2]%[ma:4c16b4d4]
% nguon: Đề tham khảo Tuyển sinh 10 TP.HCM 2025-2026 – Quận 6 – Đề 1 – Bài 2
Cho phương trình $5x^2 - 4x - 9 = 0$.
\begin{enumerate}[a)]
    \item Chứng minh phương trình trên có hai nghiệm phân biệt.
    \item Không giải phương trình, hãy tính giá trị của biểu thức $A = x_1(x_1 - 3x_2) + x_2(x_2 - 7x_1) + x_1 + x_2$.
\end{enumerate}
\loigiai{
\begin{enumerate}[a)]
    \item Phương trình $5x^2 - 4x - 9 = 0$ có $a = 5$ và $c = -9$.\\
    Vì $a \cdot c = 5 \cdot (-9) = -45 < 0$ nên phương trình luôn có hai nghiệm phân biệt $x_1, x_2$.
    \item Theo định lí Viète, ta có:
    $\left\{\begin{aligned} &S = x_1 + x_2 = \dfrac{4}{5} \\ &P = x_1 x_2 = -\dfrac{9}{5}. \end{aligned}\right.$\\
    Khai triển và biến đổi biểu thức $A$:
    $\begin{aligned}
    A &= x_1(x_1 - 3x_2) + x_2(x_2 - 7x_1) + x_1 + x_2 \\
    &= x_1^2 - 3x_1 x_2 + x_2^2 - 7x_1 x_2 + (x_1 + x_2) \\
    &= (x_1^2 + x_2^2) - 10x_1 x_2 + (x_1 + x_2) \\
    &= \left[(x_1 + x_2)^2 - 2x_1 x_2\right] - 10x_1 x_2 + (x_1 + x_2) \\
    &= S^2 - 12P + S.
    \end{aligned}$\\
    Thay các giá trị $S = \dfrac{4}{5}$ và $P = -\dfrac{9}{5}$ vào biểu thức $A$, ta được:
    $A = \left(\dfrac{4}{5}\right)^2 - 12 \cdot \left(-\dfrac{9}{5}\right) + \dfrac{4}{5} = \dfrac{16}{25} + \dfrac{108}{5} + \dfrac{4}{5} = \dfrac{16}{25} + \dfrac{112}{5} = \dfrac{16 + 560}{25} = \dfrac{576}{25} = 23{,}04$.\\
    Vậy $A = \dfrac{576}{25} = 23{,}04$.
\end{enumerate}
}
\end{bt}

% nguồn: 2026-11-15_09_de_tham_khao_quan_6_de_2_va_quan_7_de_1 (nhập 15/11/2026)
\begin{bt}%[9D6H20-2]%[ma:42d656bd]
% nguon: Đề tham khảo Tuyển sinh 10 TP.HCM 2025-2026 – Quận 6 – Đề 2 – Bài 2
Cho phương trình $5x^2 - 4x - 9 = 0$.
\begin{enumerate}[a)]
    \item Chứng minh phương trình trên có hai nghiệm phân biệt.
    \item Không giải phương trình, hãy tính giá trị của biểu thức: $B = x_1(x_1 - 3) + x_2(x_2 + 7x_1) - 3x_2$.
\end{enumerate}
\loigiai{
\begin{enumerate}[a)]
    \item Phương trình $5x^2 - 4x - 9 = 0$ có các hệ số $a = 5$ và $c = -9$.\\
    Vì $a \cdot c = 5 \cdot (-9) = -45 < 0$ nên phương trình luôn có hai nghiệm phân biệt $x_1, x_2$.
    \item Theo định lí Viète, ta có:
    $\left\{\begin{aligned} &S = x_1 + x_2 = \dfrac{4}{5} \\ &P = x_1 x_2 = -\dfrac{9}{5}. \end{aligned}\right.$\\
    Biến đổi biểu thức $B$:
    $\begin{aligned}
    B &= x_1(x_1 - 3) + x_2(x_2 + 7x_1) - 3x_2 \\
    &= x_1^2 - 3x_1 + x_2^2 + 7x_1 x_2 - 3x_2 \\
    &= (x_1^2 + x_2^2) + 7x_1 x_2 - 3(x_1 + x_2) \\
    &= \left[(x_1 + x_2)^2 - 2x_1 x_2\right] + 7x_1 x_2 - 3(x_1 + x_2) \\
    &= S^2 + 5P - 3S.
    \end{aligned}$\\
    Thay $S = \dfrac{4}{5}$ và $P = -\dfrac{9}{5}$ vào biểu thức, ta được:
    $B = \left(\dfrac{4}{5}\right)^2 + 5 \cdot \left(-\dfrac{9}{5}\right) - 3 \cdot \dfrac{4}{5} = \dfrac{16}{25} - 9 - \dfrac{12}{5} = \dfrac{16 - 225 - 60}{25} = -\dfrac{269}{25} = -10{,}76$.\\
    Vậy $B = -\dfrac{269}{25} = -10{,}76$.
\end{enumerate}
}
\end{bt}

% nguồn: 2026-11-15_09_de_tham_khao_quan_6_de_2_va_quan_7_de_1 (nhập 15/11/2026)
\begin{bt}%[9D6H20-2]%[ma:bd58128c]
% nguon: Đề tham khảo Tuyển sinh 10 TP.HCM 2025-2026 – Quận 7 – Đề 1 – Bài 2
Cho phương trình $2x^2 - 6x - 1 = 0$ có hai nghiệm phân biệt là $x_1, x_2$.
\begin{enumerate}[a)]
    \item Chứng minh phương trình trên có hai nghiệm phân biệt.
    \item Không giải phương trình, hãy tính giá trị của biểu thức: $A = \dfrac{x_1 - 2}{x_2 - 1} + \dfrac{x_2 - 2}{x_1 - 1}$.
\end{enumerate}
\loigiai{
\begin{enumerate}[a)]
    \item Phương trình $2x^2 - 6x - 1 = 0$ có $a = 2$ và $c = -1$.\\
    Vì $a \cdot c = 2 \cdot (-1) = -2 < 0$ nên phương trình luôn có hai nghiệm phân biệt $x_1, x_2$.
    \item Theo định lí Viète, ta có:
    $\left\{\begin{aligned} &S = x_1 + x_2 = \dfrac{6}{2} = 3 \\ &P = x_1 x_2 = -\dfrac{1}{2}. \end{aligned}\right.$\\
    Quy đồng mẫu thức biểu thức $A$:
    $\begin{aligned}
    A &= \dfrac{(x_1 - 2)(x_1 - 1) + (x_2 - 2)(x_2 - 1)}{(x_2 - 1)(x_1 - 1)} \\
    &= \dfrac{(x_1^2 - 3x_1 + 2) + (x_2^2 - 3x_2 + 2)}{x_1 x_2 - (x_1 + x_2) + 1} \\
    &= \dfrac{(x_1^2 + x_2^2) - 3(x_1 + x_2) + 4}{x_1 x_2 - (x_1 + x_2) + 1} \\
    &= \dfrac{\left[(x_1 + x_2)^2 - 2x_1 x_2\right] - 3(x_1 + x_2) + 4}{x_1 x_2 - (x_1 + x_2) + 1} \\
    &= \dfrac{S^2 - 2P - 3S + 4}{P - S + 1}.
    \end{aligned}$\\
    Thay $S = 3$ và $P = -\dfrac{1}{2}$ vào biểu thức:
    $\begin{aligned}
    \text{Tử số} &= 3^2 - 2 \cdot \left(-\dfrac{1}{2}\right) - 3 \cdot 3 + 4 = 9 + 1 - 9 + 4 = 5, \\
    \text{Mẫu số} &= -\dfrac{1}{2} - 3 + 1 = -\dfrac{5}{2}.
    \end{aligned}$\\
    Do đó:
    $A = \dfrac{5}{-\dfrac{5}{2}} = -2$.\\
    Vậy $A = -2$.
\end{enumerate}
}
\end{bt}
